\chapter{Relational and Path-Sensitive Analysis}
\label{ch:relational-path-analysis}

A cursor and an end pointer can both have large ranges while their distance stays small.
Two counters can grow without bound while differing by exactly three.
An optimizer interested in those differences gets little help from separate intervals, even when each interval is the best possible one.
The missing information is a relationship between values, not a better worklist order.

This chapter develops one precise relational domain, difference constraints over mathematical integers, and separates it from the independent choice of retaining path alternatives.
We use the domain to prove a guard redundant, reject an infeasible branch, and construct a loop invariant with a terminating widening policy.
The analysis order is the concretization order of Chapter~\ref{ch:designing-sound-analyses}: smaller elements admit fewer states, bottom is impossible, and join covers alternatives.

\section{A range cannot express a distance}
\label{sec:ch43-distance}

Consider an integer input $i$ with $0\leq i\leq5$, followed by $j=i+3$.
The exact reachable pairs are $(0,3),(1,4),\ldots,(5,8)$.
Separate intervals retain $i\in[0,5]$ and $j\in[3,8]$.
They also admit the spurious pair $(5,3)$, so they cannot prove $i<j$ by comparing the maximum of $i$ with the minimum of $j$.
A relation $j-i=3$ proves it immediately.

Figure~\ref{fig:ch43-relation} shows the geometric difference between the interval box and the six reachable integer points.
The relational fact can justify deleting a comparison or a bounds check whose other preconditions are already established.
It cannot by itself establish that $j$ is a valid array length, that an associated pointer is live, or that a memory access has the right element size.
Those remain separate client obligations.

\begin{figure}
\centering
\begin{tikzpicture}[x=7mm,y=5mm]
\draw[cfg edge] (-0.5,0) -- (7,0) node[right] {$i$};
\draw[cfg edge] (0,-0.5) -- (0,10) node[above] {$j$};
\draw[dashed] (0,3) rectangle (5,8);
\draw[semithick] (0,3) -- (5,8);
\foreach \x in {0,1,2,3,4,5} {
  \fill (\x,{\x+3}) circle (2pt);
}
\node[font=\small,anchor=west] at (5.6,6.6) {$j-i=3$};
\node[font=\small,anchor=west] at (5.6,3.5) {interval box};
\node[below] at (5,0) {$5$};
\node[left] at (0,3) {$3$};
\node[left] at (0,8) {$8$};
\draw (5,3) circle (3pt);
\node[font=\small,anchor=west] at (5.6,1.5) {$(5,3)$ is spurious};
\end{tikzpicture}
\caption{The six filled points satisfy the fixed difference. Independent intervals also admit points off the diagonal, including the open point $(5,3)$.}
\label{fig:ch43-relation}
\end{figure}


A \term{difference constraint} has the form $x_i-x_j\leq c$, where $c$ is an integer constant.
Introduce a distinguished variable $x_0=0$.
Then $x_i\leq u$ is $x_i-x_0\leq u$, and $x_i\geq l$ is $x_0-x_i\leq-l$.
An equality such as $j-i=3$ is represented by two constraints, $j-i\leq3$ and $i-j\leq-3$.
This domain includes intervals and fixed differences, but not general linear relations such as $x+y=8$ or $2x-y=0$.

The choice is deliberate.
We use only non-strict difference constraints with integer coefficients $+1$ and $-1$, integer constants, and an implicit zero variable.
We do not implement octagons, which additionally encode sums and require different normalization rules.
For integer variables, a strict guard $x_i<x_j+c$ is equivalent to $x_i-x_j\leq c-1$.
For real variables, that conversion would be wrong; strictness would need an explicit representation.

\section{Difference-bound matrices and closure}
\label{sec:ch43-dbm}

A \term{difference-bound matrix}, or DBM, stores one upper bound $M_{ij}$ for each difference $x_i-x_j$.
The value $+\infty$ means that this entry imposes no bound.
The concretization is the set of integer tuples satisfying every matrix entry, with $x_0=0$.
A separate bottom tag represents inconsistency.
For a nonempty state, diagonal entries are zero.
Smaller bounds are stronger constraints.

Constraints imply other constraints.
From $x_i-x_k\leq a$ and $x_k-x_j\leq b$, addition gives $x_i-x_j\leq a+b$.
A matrix is \term{closed} when every entry already contains the tightest bound obtainable by such paths.
Algorithm~\ref{alg:ch43-close} computes closure by considering each variable as an intermediate vertex.
With the convention used here, a graph edge $j\to i$ carries weight $M_{ij}$.
A negative-weight cycle asserts $x_i-x_i<0$ and proves inconsistency.

\begin{algorithm}
\caption{Close an integer DBM and detect inconsistency.}
\label{alg:ch43-close}
\begin{pseudocode}
\Require $m\times m$ DBM $M$, entries in $\mathbb{Z}\cup\{+\infty\}$
\Require index $0$ denotes zero; addition uses exact integer arithmetic
\Ensure canonical bounds for the same solutions, or $\bot$ if none exist
\State $D\gets$ a copy of $M$
\For{$i=0,\ldots,m-1$}
  \State $D_{ii}\gets\min(D_{ii},0)$
\EndFor
\For{$k=0,\ldots,m-1$}
  \For{$i=0,\ldots,m-1$}
    \For{$j=0,\ldots,m-1$}
      \If{$D_{ik}\neq+\infty$ and $D_{kj}\neq+\infty$}
        \State $D_{ij}\gets\min(D_{ij},D_{ik}+D_{kj})$
      \EndIf
    \EndFor
  \EndFor
  \If{some $D_{ii}<0$}
    \State \Return $\bot$
  \EndIf
\EndFor
\State \Return $D$
\end{pseudocode}
\end{algorithm}

For this constraint language, absence of a negative cycle implies feasibility over the integers.
Integer edge weights admit integer shortest-path potentials; translating all potentials to make $x_0=0$ preserves every difference.
Closure is therefore an exact feasibility test for a conjunction represented by one DBM.
It is not an exact feasibility test for the original program after approximation has added states.

Let $V$ be the number of tracked program variables and $m=V+1$ include zero.
The algorithm uses $O(m^3)$ bound additions and comparisons and $O(m^2)$ bound cells.
These are representation-operation counts.
If finite bounds require $b$ bits during computation, each addition and comparison also has a bit cost depending on $b$.
A machine implementation must not let host-integer overflow silently turn a large positive path length into a negative cycle.
Use exact arithmetic or a proved conservative bound representation.

For nonempty closed matrices, inclusion is entrywise comparison: $M\sqsubseteq N$ exactly when $M_{ij}\leq N_{ij}$ for every pair.
The join is entrywise maximum, which retains the weakest common bound on each difference.
For closed operands this is again a closed DBM and is the least DBM covering their union.
The meet takes entrywise minima and then closes, possibly obtaining bottom.
Bottom joins with another state to give that state and meets to bottom.
Do not substitute entrywise comparison for semantic inclusion on arbitrary noncanonical representations: one matrix can imply a tighter bound than its stored entry reveals.

\section{Transfers that preserve relations}
\label{sec:ch43-transfers}

A guard $x_a-x_b\leq c$ tightens $M_{ab}$ to $\min(M_{ab},c)$ and then closes.
If closure returns bottom, the guarded edge is infeasible for the represented state.
To forget a variable $x_a$, first close, then set its off-diagonal row and column to infinity.
Relations between the surviving variables remain; closure before projection ensures that consequences passing through $x_a$ have already been recorded.
Forgetting without that step is sound but can needlessly lose surviving relations.

Assignment $x_a=x_b+c$ with $a\neq b$ first forgets the old $x_a$, then inserts $M_{ab}\leq c$ and $M_{ba}\leq-c$, and closes.
Other variables keep their old relationships.
Assignment $x_a=x_a+c$ needs a different rule because forgetting $x_a$ would erase its relationship to the old state.
Add $c$ to each finite entry in row $a$, subtract $c$ from each finite entry in column $a$, leave $M_{aa}=0$, and leave all other entries unchanged.
This substitution is exact for a translation over mathematical integers.

For a constant assignment, use $b=0$ in the copy rule.
For an unsupported expression such as $x_a=x_b*x_c$, forget $x_a$, then optionally add sound interval bounds derived from the old state.
Retaining the old row and column after assigning a new value is unsound.
A production domain can add specialized transfers incrementally, with the forgetting transfer serving as the correct baseline.

An SSA edge containing several $\phi$-assignments is a parallel assignment.
Introduce temporary old names or apply a simultaneous substitution before projecting them away.
Serially executing a swap as two destructive copies loses the old value needed by the second copy.
The issue is semantic even when an implementation stores the matrix persistently.

\section{Compute a relational guard}
\label{sec:ch43-worked}

Use variable order $(x_0,i,j)=(0,i,j)$ for the example $0\leq i\leq5$, $j=i+3$.
Its closed matrix is
\begin{equation}
 D=\begin{pmatrix}
 0&0&-3\\
 5&0&-3\\
 8&3&0
 \end{pmatrix}.
 \label{eq:ch43-example-dbm}
\end{equation}
For example, $D_{02}=-3$ means $0-j\leq-3$, so $j\geq3$.
The entries $D_{12}=-3$ and $D_{21}=3$ encode $j-i=3$.
Every entry in Equation~\ref{eq:ch43-example-dbm} can be obtained by maximizing its difference over the six concrete pairs.

Now take a branch guarded by $j\leq6$.
It lowers $D_{20}$ from $8$ to $6$.
Closure then lowers $D_{10}$ to $D_{12}+D_{20}=-3+6=3$.
The resulting state has $0\leq i\leq3$, $3\leq j\leq6$, and still $j-i=3$.
Table~\ref{tab:ch43-bounds} compares interval refinement and relational closure.
The interval state can cap $j$ but has no recorded relationship with which to cap $i$.

\begin{table}
\centering
\caption{Closure transfers a bound on $j$ back to $i$ through their fixed difference.}
\label{tab:ch43-bounds}
\begin{tabular}{p{0.28\linewidth}p{0.27\linewidth}p{0.33\linewidth}}
\toprule
Stage & Independent intervals & Closed difference constraints \\
\midrule
After $j=i+3$ & $i\in[0,5]$, $j\in[3,8]$ & Same bounds; $j-i=3$ \\
After $j\leq6$ & $i\in[0,5]$, $j\in[3,6]$ & $i\in[0,3]$, $j\in[3,6]$; $j-i=3$ \\
Then require $i\geq4$ & $i\in[4,5]$, $j\in[3,6]$ & $\bot$ \\
\bottomrule
\end{tabular}
\end{table}

The last guard inserts $D_{01}\leq-4$ while closure has $D_{10}=3$.
The cycle through zero and $i$ has weight $-4+3=-1$, so no state remains.
Deleting this edge is justified by inconsistency of an over-approximation: if even the larger set has no state, the concrete set has none.
In contrast, a nonempty closed matrix proves only that the abstract constraints are feasible.
It does not prove that the program reaches a corresponding state.

\section{Relations and path partitions solve different losses}
\label{sec:ch43-paths}

Suppose one branch assigns $(x,y)=(1,7)$ and another assigns $(x,y)=(7,1)$.
Both establish $x+y=8$.
The interval join gives $x,y\in[1,7]$.
The DBM join adds only $-6\leq x-y\leq6$, which does not improve this box: it still admits $(1,1)$ and $(7,7)$.
Difference constraints cannot express the common sum.

A \term{path partition} retains separate abstract states for selected alternatives instead of joining them immediately.
Figure~\ref{fig:ch43-partitions} keeps the two branches until after $s=x+y$.
Each partition computes $s=8$, so their later join retains that constant even with a nonrelational value domain.
An analysis using a richer domain capable of expressing the sum could instead merge earlier.
Relational precision and path precision are independent resources.

\begin{figure}
\centering
\begin{tikzpicture}[node distance=12mm and 10mm]
\node[basic block] (entry) {branch on b};
\node[basic block,below left=of entry] (left) {x=1; y=7};
\node[basic block,below right=of entry] (right) {x=7; y=1};
\node[basic block,below=of left] (ls) {s=x+y: s=8};
\node[basic block,below=of right] (rs) {s=x+y: s=8};
\node[basic block,below=40mm of entry] (join) {join: s=8};
\draw[cfg edge] (entry) -- (left);
\draw[cfg edge] (entry) -- (right);
\draw[cfg edge] (left) -- (ls);
\draw[cfg edge] (right) -- (rs);
\draw[cfg edge] (ls) -- (join);
\draw[cfg edge] (rs) -- (join);
\end{tikzpicture}
\caption{Two analysis partitions survive the control-flow merge until the sum is evaluated. The duplicated analysis states do not require duplicated machine code.}
\label{fig:ch43-partitions}
\end{figure}


Formally, a finite family $\{a_1,\ldots,a_k\}$ denotes $\bigcup_i\gamma(a_i)$.
An edge transfer applies to every member; a branch can split members further.
To bound memory, an implementation can merge selected members by the underlying join.
That replacement enlarges the denotation and is sound.
Dropping a member merely because its path is unlikely is not sound for optimization.
When the partition budget is exhausted, merge or fall back to an unknown cover.

Stable partition keys matter in loops.
Partitioning by the entire branch history creates indefinitely many keys.
A finite set of predicates, a bounded loop-iteration prefix, or selected branch identifiers supplies a finite policy.
Even with finitely many keys, each partition's numerical state can have an infinite ascending chain and need widening.
Budgeting partitions does not solve numerical convergence.

The compiler should partition where a specific query benefits.
For a bounds check, the relevant distinction may be whether an earlier length test succeeded.
Retaining an unrelated logging flag doubles states without strengthening the proof.
Liveness-based projection can remove dead numerical variables, while a relevance slice can identify predicates that influence the checked address or index.
These are precision heuristics; their sound fallback remains merging.

\section{Symbolic execution and feasibility answers}
\label{sec:ch43-symbolic}

\term{Symbolic execution} records expressions over symbolic inputs together with a path condition.
For the two branches above, a symbolic state can keep $x=1,y=7$ under condition $b$, and $x=7,y=1$ under $\neg b$.
To prove an assertion $P$ on one symbolic path, ask whether its path condition together with $\neg P$ is satisfiable.
An unsatisfiable answer excludes a violating execution on that path under the encoded semantics.
The proof covers the program only if every relevant path is covered as well.

A solver answer of unknown, a timeout, or an unsupported operation is not an infeasibility proof.
Keep the path, merge it into a conservative state, or decline the optimization.
A satisfiable model establishes a concrete counterexample only if the symbolic encoding faithfully models the execution, including memory and exceptional outcomes.
If the encoding over-approximates an unknown call, its satisfying model may be another abstract witness requiring concrete validation.

DBM closure is a specialized solver for one conjunction of difference constraints.
It has predictable polynomial representation cost and supports join directly.
General symbolic conditions can retain disjunctions, multiplication, bit-vector operations, and heap encodings, but their solving cost and termination behavior differ.
Neither representation automatically dominates the other for a compiler workload.
One useful architecture computes a cheap invariant first, then sends selected failed queries to a stronger solver with that invariant as an additional assumption.
The invariant must be valid at the queried program point and revision.

Bounded loop exploration is useful for finding bugs but insufficient for proving an unbounded loop safe.
An unwinding assertion can check that the explored bound covers all executions; alternatively, an inductive invariant closes the unexplored suffix.
The next section constructs such an invariant without enumerating loop trips.

\section{Widen a representation without restoring discarded constraints}
\label{sec:ch43-widening}

Consider $i=0$, $j=3$, followed by a loop that tests $i<6$ and increments both variables by one.
At the header, the exact relation is $j-i=3$ with $0\leq i\leq6$.
The upper endpoint grows through $0,1,2,\ldots,6$ during naive ascending iteration.
Changing six to a large constant makes iteration expensive even though the relational shape remains simple.

For a loop with a loop-free body, let $E$ be its closed entry DBM and let $H(D)$ be the closed join of $E$ with the body result from the true guard applied to $D$.
Assume $H$ is sound and monotone.
The header invariant must satisfy $H(D)\sqsubseteq D$.
This formulation includes the entry contribution on every evaluation, not only during initialization.

Define a representation-level widening on matrices by retaining an old entry when the new bound is no larger and replacing it by infinity otherwise.
The difficulty is that closure can rederive an entry previously discarded by widening.
Applying closure to the left-hand accumulator before every widening invalidates the simple argument that each entry is discarded at most once.
Algorithm~\ref{alg:ch43-widen-loop} therefore keeps an unclosed accumulator $W$ and closes only a copy for semantic transfer and final queries.
It never copies those rederived bounds back into $W$.

\begin{algorithm}
\caption{Solve one DBM loop header with raw widening and bounded replay.}
\label{alg:ch43-widen-loop}
\begin{pseudocode}
\Require nonempty closed entry $E$, monotone sound header map $H$
\Require $H$ includes $E$ and returns a closed DBM; replay budget $q\geq0$
\Ensure closed header invariant $D$ with $H(D)\sqsubseteq D$
\State $W\gets E$
\Loop
  \State $C\gets H(\operatorname{Close}(W))$
  \If{$C_{ij}\leq W_{ij}$ for every $i,j$}
    \State \Break
  \EndIf
  \ForAll{pairs $(i,j)$}
    \If{$C_{ij}>W_{ij}$}
      \State $W_{ij}\gets+\infty$
    \EndIf
  \EndFor
\EndLoop
\State $D\gets\operatorname{Close}(W)$
\For{$t=1,\ldots,q$}
  \State $D\gets H(D)$
\EndFor
\State \Return $D$
\end{pseudocode}
\end{algorithm}

The accumulator only loses constraints, so it remains consistent and covers $E$.
Each changing iteration discards at least one of its at most $m(m-1)$ finite off-diagonal entries.
Thus the widening loop performs at most $m(m-1)+1$ evaluations, including the final unchanged one.
At exit, every constraint stored in $W$ is respected by $H(\operatorname{Close}(W))$, so closure of $W$ is an invariant.
The entrywise stopping test is sufficient even when $W$ is not canonical; it may be stricter than necessary.

If one evaluation of $H$ costs $T_H$ representation operations, the total bound is $O(m^2(m^3+T_H)+qT_H)$, including repeated closure of the accumulator.
This deliberately loose bound belongs to the displayed single-header algorithm.
It is not a bound for an arbitrary nested-CFG solver; nested and irreducible cycles need a separately specified iteration strategy.
Peak storage includes $O(m^2)$ matrix cells plus the body's temporary analysis states.

The replay phase is a simple bounded narrowing strategy.
From $H(D)\sqsubseteq D$, monotonicity gives $H(H(D))\sqsubseteq H(D)$.
Consequently each replay remains inductive while becoming no larger.
No claim of termination for unbounded descending iteration is needed because $q$ is finite.
Intersecting with a guessed bound, without this argument or a fresh invariant check, would not be justified.

\begin{table}
\centering
\caption{Immediate widening keeps the fixed difference, and one replay restores the finite loop bound.}
\label{tab:ch43-loop}
\begin{tabular}{p{0.25\linewidth}p{0.22\linewidth}p{0.22\linewidth}p{0.19\linewidth}}
\toprule
Stage & Range of $i$ & Range of $j$ & Relation \\
\midrule
Entry $E$ & $[0,0]$ & $[3,3]$ & $j-i=3$ \\
First $H(E)$ & $[0,1]$ & $[3,4]$ & $j-i=3$ \\
Widened $W$ & $[0,+\infty)$ & $[3,+\infty)$ & $j-i=3$ \\
$H(\operatorname{Close}(W))$ & $[0,6]$ & $[3,9]$ & $j-i=3$ \\
After one replay & $[0,6]$ & $[3,9]$ & $j-i=3$ \\
Exit guard $i\geq6$ & $[6,6]$ & $[9,9]$ & $j-i=3$ \\
\bottomrule
\end{tabular}
\end{table}

Table~\ref{tab:ch43-loop} computes the result.
The first widening drops only the growing upper bounds on $i$ and $j$; both difference entries stay fixed at $3$ and $-3$.
From the widened state, the true guard gives $i\leq5$, and the relation gives $j\leq8$.
The increments produce $i\leq6,j\leq9$, and joining entry gives the header bound used in the replay.
The exit guard then forces $i=6$, hence $j=9$.
The result is stronger than separate intervals at exit because it retains the equality linking the counters.

\section{Choose the arithmetic and the precision budget explicitly}
\label{sec:ch43-engineering}

All transfers above use mathematical integers.
For wrapping arithmetic, incrementing a counter can destroy a difference even if both counters wrap individually.
For example, unsigned eight-bit values $i=252,j=255$ have mathematical difference three; incrementing both modulo $256$ gives $i=253,j=0$.
A compiler can apply the translation transfer only after proving no relevant wrap occurs, or after adopting a domain whose relations are defined modulo the bit width.
The latter is a different domain with different guard rules.

Our counted-loop example stays between zero and nine, so it is also valid for machine integer types that represent those values and perform the displayed operations without exceptional outcomes.
That justification comes from the established bounds, not from assuming that all loops behave like mathematical integers.
Pointer differences additionally require the language's object and provenance rules; numerical closeness alone does not establish a legal pointer subtraction.

Dense DBMs become expensive as the number of live variables grows.
Track selected packs of related variables rather than one matrix for every scalar in a large function.
Projection at lifetime boundaries removes dead names, while overlap between packs can exchange interval bounds.
Any exchange must preserve the joint concretization; a reduction may strengthen one component only using facts proved by the others.
Scheduling fewer closures is another engineering option, provided every operation that requires canonical input actually receives it.

Precision should be measured at clients.
Count checks proved redundant, queries left unknown, matrix dimensions, closure operations, and partitions merged.
A more precise invariant that enables no additional transformation may not justify its compile-time cost.
Conversely, a small targeted relational pack can pay for itself by proving one hot-loop bounds check redundant.
The soundness contract remains the same when the budget changes: exhausted precision yields a larger cover, never discarded behavior.

\section{Exercises}
\label{sec:ch43-exercises}

\begin{enumerate}
\item \label{ex:ch43-close}
Let $0\leq x\leq4$ and $y=x+2$.
Write the closed DBM in order $(0,x,y)$.
Add $y\leq4$, close it, and identify every changed entry.
\item \label{ex:ch43-cycle}
For that refined state, add $x\geq3$.
Give a negative cycle and its total weight.
Explain why a nonempty DBM would not by itself establish a reachable program path.
\item \label{ex:ch43-partition}
Two arms assign $(x,y)=(2,9)$ and $(9,2)$.
Compute $s=x+y$ after a DBM join and after separate per-arm evaluation.
Give a pair admitted by the joined DBM that explains the precision loss.
\item \label{ex:ch43-loop}
Change the loop to $i=0,j=4$, guard $i<7$, and increments by one.
Run immediate widening and one replay as in Algorithm~\ref{alg:ch43-widen-loop}.
Give the widened, replayed, and exit states, including the relation.
\end{enumerate}

\section{Further reading}
\label{sec:ch43-reading}

Cousot and Cousot explain why widening and inductive over-approximations can replace exact iteration without sacrificing safety~\cite{cousot-1977-abstract}.
Nielson, Nielson, and Hankin connect domain design, abstract transfer functions, and fixed-point construction; use their framework to check which precision claims depend on the chosen representation~\cite{nielson-1999-principles}.
Baier and Katoen provide complementary foundations for state-space exploration and the interpretation of counterexamples~\cite{baier-2008-model-checking}.
